\documentclass[a4paper,12pt]{article}
\usepackage{color}
\usepackage{graphicx, gensymb}
\usepackage{amsfonts, cancel, amssymb}
\usepackage{amsmath, amsthm, array, esvect, esint}
\usepackage{typearea}
\usepackage[a4paper,left=2cm,right=2cm,top=2cm,bottom=3cm,bindingoffset=5mm]{geometry}
\newcommand\numberthis{\addtocounter{equation}{1}\tag{\theequation}}
\newcommand{\bra}{}
\newcommand{\kom}{}
\newcommand{\brak}{}
\newcommand{\braket}{}
\newcommand{\intx}{}
\newcommand{\intr}{}
\newcommand{\kla}{}
\newcommand{\klam}{}
\newcommand{\klammer}{}
\newcommand{\oper}{}
\newcommand{\tecxt}{}

\begin{document}

\title{08 Auswahlregeln}
\author{Joachim Schwardt}
\date{\today}
\maketitle

\section*{Aufgabe 20: Feinstruktur und Auswahlregeln}
\section*{Aufgabe 21: Symmetrie und Auswahlregeln}
Sei $H=H(r,\vartheta)$. Dann gilt mit $\phi=:g(r,\vartheta)f(\varphi)$:
\begin{align*}
E\phi&=H\phi=-\frac{\hbar^2}{2m}\nabla^2gf +V(r,\vartheta)gf\ \ \Rightarrow\ \ \frac{2mr^2}{\hbar^2}(E-V(r,\vartheta))+r^2\frac{\nabla_{r,\vartheta}^2g}{g}=-\frac{1}{\sin^2\varphi}\frac{f''}{f}\\
&\Rightarrow\ \ 0=k\sin^2\varphi f+f''\\
L_z\phi&=\hbar \tecxt\text{m}\phi\ \ \Rightarrow\ \ \frac{\hbar}{\tecxt\text{i}}f'=\hbar \tecxt\text{m} f\ \ \Rightarrow\ \ f=e^{\tecxt\text{im}\varphi}\\
f(2\pi)&\overset{!}{=}f(0)\ \ \Rightarrow\ \ \tecxt\text{m}\in\mathbb{Z}
\end{align*}
\section*{Aufgabe 22: Atomare Übergänge}
\subsection*{a)}
Es gilt $N'(t)=-\frac{\Gamma}{\hbar}N(t)$. Daraus folgt direkt $N(t)=N_0e^{-\frac{\Gamma}{\hbar}t}$.
\subsection*{b)}
Hier ist $N'(t)=-\frac{\Gamma_1}{\hbar}N(t)-\frac{\Gamma_2}{\hbar}N(t)$. Es gibt also die gleiche Lösung, $N(t)=N_0e^{-\frac{\Gamma_1+\Gamma_2}{\hbar}t}$. Die Lebensdauer ist somit $T=\frac{\hbar}{\Gamma_1+\Gamma_2}$.
\subsection*{c)}





































\newpage
1
\end{document}